\section{System Design}
\label{sec:design}

Pigs is structured as a three-layer system: a \emph{proxy} layer that routes HTTP requests, a \emph{phased runtime} layer that orchestrates the Pre$\rightarrow$Executor$\rightarrow$Post state machine, and a \emph{CLI} layer that provides an interactive agent interface. This section describes the core design of each layer.
\begin{chinese}
Pigs 被组织为一个三层系统：负责路由 HTTP 请求的\emph{代理}层、负责编排 Pre$\rightarrow$Executor$\rightarrow$Post 状态机的\emph{分阶段运行时}层，以及提供交互式智能体界面的 \emph{CLI} 层。本节描述各层的核心设计。
\end{chinese}

\subsection{Architecture Overview}
\label{sec:overview}
\begin{chinese}
\subsection{架构概览}
\end{chinese}

The system operates as a single-process server listening on one port (default 3927). All three API protocols are served from the same port, distinguished by request path:
\begin{chinese}
系统作为单进程服务器运行，监听一个端口（默认 3927）。三种 API 协议均从同一端口提供服务，通过请求路径区分：
\end{chinese}

\begin{itemize}[nosep]
    \item \texttt{/chat/completions} $\rightarrow$ OpenAI Chat
    \item \texttt{/v1/messages} $\rightarrow$ Anthropic Messages
    \item \texttt{/responses} $\rightarrow$ OpenAI Responses
\end{itemize}

Requests are routed by model identifier: models without the \texttt{-pig} suffix are passed through to the upstream provider (with retry, body cleaning, and thinking-effort injection); models with the \texttt{-pig} suffix enter the phased runtime. Each configured model is automatically exposed as two identifiers: \texttt{\{model\}} (passthrough) and \texttt{\{model\}-pig} (phased).
\begin{chinese}
请求通过模型标识符路由：不带 \texttt{-pig} 后缀的模型直接透传至上游提供商（附带重试、主体清洗和思考力度注入）；带 \texttt{-pig} 后缀的模型进入分阶段运行时。每个配置的模型自动暴露为两个标识符：\texttt{\{model\}}（透传）和 \texttt{\{model\}-pig}（分阶段）。
\end{chinese}

\Cref{fig:arch} illustrates the request flow. When a \texttt{-pig} request arrives, the proxy constructs a \texttt{HttpRequestEnvelope} and hands it to the \texttt{HttpPhasedRuntime}. Each phase subrequest is sent back to the proxy via HTTP loopback (with an internal token header), re-entering the passthrough path with the \texttt{-pig} suffix removed. This design ensures that phase subrequests transparently inherit all proxy-level features.
\begin{chinese}
\Cref{fig:arch} 展示了请求流。当 \texttt{-pig} 请求到达时，代理构建一个 \texttt{HttpRequestEnvelope} 并将其交给 \texttt{HttpPhasedRuntime}。每个阶段子请求通过 HTTP 回环（附带一个内部令牌头部）发送回代理，在去除 \texttt{-pig} 后缀后重新进入透传路径。这一设计确保阶段子请求透明地继承所有代理级特性。
\end{chinese}

\begin{figure}[t]
\centering
\begin{tikzpicture}[
  node distance=0.8cm,
  box/.style={draw, rounded corners, minimum width=3.5cm, minimum height=0.8cm, align=center, font=\small},
  client/.style={box, fill=blue!10},
  proxy/.style={box, fill=green!10},
  runtime/.style={box, fill=orange!10},
  provider/.style={box, fill=red!10},
  arr/.style={-Stealth, thick},
  loopback/.style={-Stealth, thick, dashed, draw=purple!70},
]
\node[client] (client) {Client (model: X-pig)};
\node[proxy, below=of client] (proxy1) {pigs-proxy\\(route by -pig)};
\node[runtime, below=of proxy1] (runtime) {HttpPhasedRuntime\\Pre $\to$ Executor $\to$ Post};
\node[proxy, below=of runtime] (proxy2) {pigs-proxy\\(passthrough, model: X)};
\node[provider, below=of proxy2] (provider) {Upstream Provider};

\draw[arr] (client) -- (proxy1);
\draw[arr] (proxy1) -- (runtime);
\draw[loopback] (runtime) -- node[right, font=\scriptsize, text=purple!70] {HTTP loopback} (proxy2);
\draw[arr] (proxy2) -- (provider);
\end{tikzpicture}
\caption{Request flow for a \texttt{-pig} model. Phase subrequests re-enter the proxy via HTTP loopback (dashed), inheriting passthrough features.}
\label{fig:arch}
\end{figure}
\subsection{Phase State Machine}
\label{sec:state}
\begin{chinese}
\subsection{阶段状态机}
\end{chinese}

The core of Pigs is a pure state machine with three phases and three transition signals. The state is encapsulated in \texttt{OrchestrationState}, which tracks the current phase, accumulated outputs, failure history, and iteration counters.
\begin{chinese}
Pigs 的核心是一个具有三个阶段和三个转移信号的纯状态机。状态封装在 \texttt{OrchestrationState} 中，用于跟踪当前阶段、累积输出、失败历史和迭代计数器。
\end{chinese}

\begin{algorithm}[t]
\caption{Phase state machine transition}
\label{alg:state}
\begin{algorithmic}[1]
\Function{advance}{output}
  \State $phase \gets \text{current phase}$
  \If{$phase = \text{Pre}$}
    \If{$\text{detect}(output) = \texttt{PIGEND}$}
      \State \Return Complete \Comment{Short-circuit: trivially simple task}
    \ElsIf{$\text{detect}(output) = \texttt{PIGFAIL}$}
      \State \Return replan(output) \Comment{Failure: return to Pre}
    \Else
      \State $\text{pre\_output} \gets \text{strip}(output)$
      \State \Return Continue(Executor)
    \EndIf
  \ElsIf{$phase = \text{Executor}$}
    \State $\text{executor\_outputs} \gets \text{strip}(output)$
    \State \Return Continue(Post) \Comment{Always proceed to Post}
  \Else \Comment{$phase = \text{Post}$}
    \If{$\text{detect}(output) = \texttt{PIGEND}$}
      \State \Return Complete
    \ElsIf{$\text{detect}(output) = \texttt{PIGFAIL}$}
      \State \Return replan(output)
    \Else
      \If{$\text{post\_iterations} \geq \text{max}$}
        \State \Return Error(PostBudgetExceeded)
      \EndIf
      \State $\text{post\_iterations} \mathrel{+}= 1$
      \State \Return Continue(Post) \Comment{Stay in Post}
    \EndIf
  \EndIf
\EndFunction
\end{algorithmic}
\end{algorithm}

The state machine has several key properties:
\begin{chinese}
该状态机具有若干关键特性：
\end{chinese}

\textbf{Markerless Post stays in Post.} When the Post phase produces output without a control marker, the state machine continues in Post (not back to Executor). This design choice prevents infinite Executor$\leftrightarrow$Post loops: the model is expected to make incremental progress within Post until it either completes or fails.
\begin{chinese}
\textbf{无标记的 Post 停留在 Post。} 当 Post 阶段产生不含控制标记的输出时，状态机继续停留在 Post 阶段（而非返回 Executor）。这一设计选择防止了 Executor$\leftrightarrow$Post 的无限循环：模型应在 Post 内部逐步推进，直至完成或失败。
\end{chinese}

\textbf{Budget exhaustion is an error, not success.} When the Post iteration budget is exhausted, the state machine returns an \texttt{OrchestrationError}, not a successful completion. This prevents the system from silently accepting incomplete work.
\begin{chinese}
\textbf{预算耗尽即为错误，而非成功。} 当 Post 迭代预算耗尽时，状态机返回 \texttt{OrchestrationError}，而非成功完成。这防止了系统静默接受不完整的工作。
\end{chinese}

\textbf{Failure preserves history.} When \texttt{PIGFAIL} triggers a replan, the failure output is preserved and presented to the next Pre phase as numbered failure records (``1st failure: \ldots'', ``2nd failure: \ldots''). This allows the model to learn from previous attempts.
\begin{chinese}
\textbf{失败保留历史。} 当 \texttt{PIGFAIL} 触发重新规划时，失败输出会被保留并以编号的失败记录形式呈现给下一个 Pre 阶段（"第1次失败：\ldots"、"第2次失败：\ldots"）。这使得模型能够从先前的尝试中学习。
\end{chinese}

\begin{figure}[t]
\centering
\begin{tikzpicture}[
  node distance=1.8cm and 2.2cm,
  state/.style={draw, circle, minimum size=1.4cm, align=center, font=\small},
  terminal/.style={draw, rounded corners, minimum width=1.6cm, minimum height=0.7cm, align=center, font=\small, fill=gray!10},
  arr/.style={-Stealth, thick},
  fail/.style={-Stealth, thick, dashed, draw=red!70},
  every label/.style={font=\scriptsize},
]
\node[state, fill=blue!10] (pre) {Pre};
\node[state, fill=orange!10, right=of pre] (exec) {Exec.};
\node[state, fill=green!10, right=of exec] (post) {Post};
\node[terminal, right=of post, fill=green!20] (done) {Complete};
\node[terminal, below=1.6cm of post, fill=red!10] (error) {Error};

\draw[arr] (pre) -- node[above, font=\scriptsize] {no marker} (exec);
\draw[arr] (exec) -- node[above, font=\scriptsize] {always} (post);
\draw[arr] (post) -- node[above, font=\scriptsize] {\texttt{PIGEND}} (done);
\draw[arr] (pre.north) .. controls +(0, 1.2) and +(0, 1.2) .. node[above, font=\scriptsize] {\texttt{PIGEND} (short-circuit)} (done.north);
\draw[arr] (post.east) .. controls +(0.8, -0.5) and +(0.8, 0.5) .. node[right, font=\scriptsize] {no marker} (post.south east);
\draw[fail] (post.south west) .. controls +(-1.5, -1.2) and +(1.5, -1.2) .. node[below, font=\scriptsize, text=red!70] {\texttt{PIGFAIL} (replan)} (pre.south);
\draw[fail] (post) -- node[right, font=\scriptsize, text=red!70] {budget} (error);
\end{tikzpicture}
\caption{Phase state machine transitions. Solid arrows are forward transitions; dashed red arrows are failure and budget-exhaustion paths. The curved arrow from Post back to Pre represents PIGFAIL-triggered replanning. The self-loop on Post (right side) represents markerless continuation.}
\label{fig:state}
\end{figure}

\Cref{fig:state} illustrates the complete state machine. The Pre phase can short-circuit to Complete (trivially simple tasks), and the Post phase has a self-loop for markerless continuation.
\begin{chinese}
\Cref{fig:state} 展示了完整的状态机。Pre 阶段可短路至 Complete（极其简单的任务），Post 阶段具有用于无标记继续的自环。
\end{chinese}

\subsection{Control Markers}
\label{sec:markers}
\begin{chinese}
\subsection{控制标记}
\end{chinese}

Control markers are short text tokens that the model appends to its output to signal phase transitions:
\begin{chinese}
控制标记是模型附加到其输出中的简短文本标记，用于发出阶段转移信号：
\end{chinese}

\begin{itemize}[nosep]
    \item \texttt{PIGEND}: signals successful completion. The turn ends, and the accumulated visible text is returned to the client.
    \begin{chinese}
    \texttt{PIGEND}：表示成功完成。本轮结束，累积的可见文本返回给客户端。
    \end{chinese}
    \item \texttt{PIGFAIL}: signals that the current approach has failed and the task should be replanned from Pre.
    \begin{chinese}
    \texttt{PIGFAIL}：表示当前方法已失败，任务应从 Pre 重新规划。
    \end{chinese}
\end{itemize}

Marker detection follows two rules:
\begin{chinese}
标记检测遵循两条规则：
\end{chinese}

\textbf{Last non-empty line only.} Only the final non-empty line of the model's output is checked for a marker. Earlier occurrences (e.g., the marker appearing in a sentence like ``the system outputs PIGEND when done'') are ignored. This prevents false positives from prose mentions.
\begin{chinese}
\textbf{仅检测最后一个非空行。} 仅检查模型输出的最后一个非空行是否包含标记。早期出现的标记（例如标记出现在"系统在完成时输出 PIGEND"这样的句子中）会被忽略。这防止了散文中提及标记导致的误报。
\end{chinese}

\textbf{Reason required.} A marker is valid only if at least one non-marker line precedes it. A bare \texttt{PIGEND} with no preceding content is rejected, forcing the model to provide a reason before signaling completion.
\begin{chinese}
\textbf{要求附带理由。} 标记仅在前面至少有一个非标记行时才有效。单独的 \texttt{PIGEND} 且无前置内容将被拒绝，从而强制模型在发出完成信号前提供理由。
\end{chinese}

Markers are stripped from the visible output via \texttt{strip\_markers}, which removes exact control lines while preserving prose that happens to contain marker words.
\begin{chinese}
标记通过 \texttt{strip\_markers} 从可见输出中剥离，该函数移除精确的控制行，同时保留恰好包含标记词的散文。
\end{chinese}

\subsection{Protocol-Native Request Envelope}
\label{sec:envelope}
\begin{chinese}
\subsection{协议原生请求封装}
\end{chinese}

The \texttt{HttpRequestEnvelope} is the central data structure that preserves the complete protocol-native request:
\begin{chinese}
\texttt{HttpRequestEnvelope} 是保留完整协议原生请求的核心数据结构：
\end{chinese}

\begin{lstlisting}
pub struct HttpRequestEnvelope {
    pub method: String,           // HTTP method
    pub path_and_query: String,   // path + query string
    pub headers: Vec<HeaderPair>, // raw header bytes
    pub body: Value,              // complete JSON body
    pub protocol: Protocol,       // Chat / Anthropic / Responses
    pub client_model: String,     // e.g. "gpt-4-pig"
    pub real_model: String,       // e.g. "gpt-4" (-pig stripped)
    pub stream: bool,             // streaming requested?
    // ... internal fields for current user location,
    //     tool result IDs, tool result groups
}
\end{lstlisting}

Each phase clones the envelope and applies \emph{point mutations}:
\begin{chinese}
每个阶段克隆该封装并施加\emph{定点修改}：
\end{chinese}

\begin{enumerate}[nosep]
    \item \textbf{Model}: replaced with the real (non-\texttt{-pig}) model name.
    \begin{chinese}
    \textbf{模型}：替换为真实（非 \texttt{-pig}）模型名。
    \end{chinese}
    \item \textbf{Current user input}: appended with the phase-specific prompt suffix (Pre questions, Executor plan, Post review instructions).
    \begin{chinese}
    \textbf{当前用户输入}：附加阶段特定的提示后缀（Pre 问题、Executor 计划、Post 审查指令）。
    \end{chinese}
    \item \textbf{Phase transcript}: native assistant/tool messages from prior phases are appended to the message history.
    \begin{chinese}
    \textbf{阶段记录}：来自先前阶段的原生 assistant/tool 消息被附加到消息历史中。
    \end{chinese}
    \item \textbf{Stream}: set to \texttt{true} for streaming phase subrequests.
    \begin{chinese}
    \textbf{流式}：对于流式阶段子请求设为 \texttt{true}。
    \end{chinese}
\end{enumerate}

All other fields---system/instructions, historical messages, tools/tool\_choice, reasoning parameters, cache controls, metadata, media blocks, and unknown provider extensions---remain untouched. This ensures that provider-specific features (e.g., Anthropic's \texttt{cache\_control}, OpenAI's \texttt{reasoning\_effort}) pass through without modification.
\begin{chinese}
所有其他字段——system/instructions、历史消息、tools/tool\_choice、推理参数、缓存控制、元数据、媒体块以及未知的提供商扩展——均保持不变。这确保了提供商特定特性（如 Anthropic 的 \texttt{cache\_control}、OpenAI 的 \texttt{reasoning\_effort}）原样透传，无需修改。
\end{chinese}

\subsection{HTTP Loopback Transport}
\label{sec:loopback}
\begin{chinese}
\subsection{HTTP 回环传输}
\end{chinese}

Each phase subrequest is sent back to the local proxy via HTTP, not through an in-process function call. The \texttt{LoopbackPhaseTransport} implements the \texttt{PhaseTransport} trait:
\begin{chinese}
每个阶段子请求通过 HTTP 发送回本地代理，而非通过进程内函数调用。\texttt{LoopbackPhaseTransport} 实现了 \texttt{PhaseTransport} trait：
\end{chinese}

\begin{itemize}[nosep]
    \item Converts wildcard listen addresses (\texttt{0.0.0.0}, \texttt{[::]}) to loopback (\texttt{127.0.0.1}, \texttt{::1}).
    \begin{chinese}
    将通配监听地址（\texttt{0.0.0.0}、\texttt{[::]}）转换为回环地址（\texttt{127.0.0.1}、\texttt{::1}）。
    \end{chinese}
    \item Disables environment HTTP proxy settings (\texttt{.no\_proxy()}) to prevent external proxy interception.
    \begin{chinese}
    禁用环境 HTTP 代理设置（\texttt{.no\_proxy()}），以防止外部代理拦截。
    \end{chinese}
    \item Adds a random internal token header (\texttt{x-pigs-phase-loopback}) that the proxy validates to identify legitimate phase subrequests.
    \begin{chinese}
    添加随机内部令牌头部（\texttt{x-pigs-phase-loopback}），由代理验证以识别合法的阶段子请求。
    \end{chinese}
    \item Strips hop-by-hop headers and the internal token before forwarding to the upstream provider.
    \begin{chinese}
    在转发至上游提供商前剥离逐跳头部和内部令牌。
    \end{chinese}
\end{itemize}

This design has two benefits. First, phase subrequests transparently inherit all proxy-level features: channel selection, model mapping, body cleaning (e.g., removing empty content messages), thinking-effort injection, and retry logic. Second, the control flow is explicit and debuggable---each phase subrequest appears in proxy logs as a normal HTTP request, and the internal token distinguishes it from client requests.
\begin{chinese}
此设计有两个优势。第一，阶段子请求透明地继承所有代理级特性：通道选择、模型映射、主体清洗（如移除空内容消息）、思考力度注入和重试逻辑。第二，控制流是显式且可调试的——每个阶段子请求在代理日志中作为普通 HTTP 请求出现，内部令牌使其与客户端请求区分开来。
\end{chinese}

\subsection{Bounded In-Memory Continuation}
\label{sec:continuation}
\begin{chinese}
\subsection{有界内存续传}
\end{chinese}

When the model returns tool calls during a phase, the runtime pauses and returns the tool calls to the client (or external agent). The complete phase state---including the orchestration state, phase transcript, visible text, and pending tool IDs---is stored in a \texttt{ContinuationStore}:
\begin{chinese}
当模型在某个阶段返回工具调用时，运行时暂停并将工具调用返回给客户端（或外部智能体）。完整的阶段状态——包括编排状态、阶段记录、可见文本和待处理工具 ID——存储在 \texttt{ContinuationStore} 中：
\end{chinese}

\begin{itemize}[nosep]
    \item \textbf{TTL}: each continuation expires after a configurable duration (default 30 minutes).
    \begin{chinese}
    \textbf{TTL}：每个续传在可配置的时长后过期（默认 30 分钟）。
    \end{chinese}
    \item \textbf{Capacity}: at most $N$ continuations are retained (default 256); the oldest is evicted when capacity is reached.
    \begin{chinese}
    \textbf{容量}：最多保留 $N$ 个续传（默认 256）；达到容量时驱逐最旧的续传。
    \end{chinese}
    \item \textbf{Tombstones}: when a continuation is expired, evicted, or consumed, a tombstone is placed for each tool ID, enabling distinct error reporting (Expired vs.\ Evicted vs.\ Consumed vs.\ Unknown).
    \begin{chinese}
    \textbf{墓碑}：当续传过期、被驱逐或被消费时，为每个工具 ID 放置一个墓碑，实现差异化错误报告（Expired 与 Evicted 与 Consumed 与 Unknown）。
    \end{chinese}
    \item \textbf{No auth headers}: authentication headers are cleared before storing; the continuation uses the incoming request's headers on resume.
    \begin{chinese}
    \textbf{不存储认证头部}：认证头部在存储前被清除；续传在恢复时使用传入请求的头部。
    \end{chinese}
\end{itemize}

When the next request arrives with tool results, the runtime matches them by tool-call ID, validates protocol and model consistency, and resumes the paused phase. Partial parallel results (multiple tool calls returning in separate requests) are accumulated until all pending IDs are present.
\begin{chinese}
当下一个请求携带工具结果到达时，运行时通过工具调用 ID 匹配它们，验证协议和模型一致性，并恢复暂停的阶段。部分并行结果（多个工具调用在单独请求中返回）会被累积，直到所有待处理 ID 均到齐。
\end{chinese}

\subsection{Continuous SSE Output}
\label{sec:sse}
\begin{chinese}
\subsection{连续 SSE 输出}
\end{chinese}

For streaming requests, the runtime produces a continuous SSE stream that concatenates visible text from all phases. A \emph{marker line buffer} holds back the last non-empty line of each phase's output until the phase completes, then checks for control markers. If a marker is found, the held-back line is stripped; otherwise, it is forwarded. This prevents \texttt{PIGEND}/\texttt{PIGFAIL} from leaking into the client-visible stream, even when markers are split across network chunks.
\begin{chinese}
对于流式请求，运行时产生连续的 SSE 流，拼接所有阶段的可见文本。\emph{标记行缓冲区}暂缓每个阶段输出的最后一个非空行，直到该阶段完成，然后检查控制标记。若发现标记，被暂缓的行将被剥离；否则，将其转发。这防止了 \texttt{PIGEND}/\texttt{PIGFAIL} 泄漏到客户端可见的流中，即使标记跨越网络数据块的分界也是如此。
\end{chinese}

The outer SSE encoder generates protocol-native event sequences for all three protocols:
\begin{chinese}
外层 SSE 编码器为所有三种协议生成协议原生的事件序列：
\end{chinese}
\begin{itemize}[nosep]
    \item \textbf{Chat}: \texttt{chat.completion.chunk} deltas with role/content/tool\_calls, terminated by \texttt{finish\_reason} + \texttt{[DONE]}.
    \begin{chinese}
    \textbf{Chat}：\texttt{chat.completion.chunk} 增量，附带 role/content/tool\_calls，以 \texttt{finish\_reason} + \texttt{[DONE]} 终止。
    \end{chinese}
    \item \textbf{Anthropic}: \texttt{message\_start} $\rightarrow$ \texttt{content\_block\_start/delta/stop} $\rightarrow$ \texttt{message\_delta} $\rightarrow$ \texttt{message\_stop}.
    \begin{chinese}
    \textbf{Anthropic}：\texttt{message\_start} $\rightarrow$ \texttt{content\_block\_start/delta/stop} $\rightarrow$ \texttt{message\_delta} $\rightarrow$ \texttt{message\_stop}。
    \end{chinese}
    \item \textbf{Responses}: \texttt{response.created} $\rightarrow$ \texttt{output\_item/content\_part/text\_delta/done} $\rightarrow$ \texttt{response.completed}, with monotonic sequence numbers.
    \begin{chinese}
    \textbf{Responses}：\texttt{response.created} $\rightarrow$ \texttt{output\_item/content\_part/text\_delta/done} $\rightarrow$ \texttt{response.completed}，带有单调递增的序列号。
    \end{chinese}
\end{itemize}

Error streams emit protocol-native error events (Chat: \texttt{error} in data; Anthropic: \texttt{event: error}; Responses: \texttt{response.failed}) without any success terminal frame.
\begin{chinese}
错误流发出协议原生的错误事件（Chat：数据中的 \texttt{error}；Anthropic：\texttt{event: error}；Responses：\texttt{response.failed}），不含任何成功的终止帧。
\end{chinese}
